\section{Критерий согласия для \(X_4\)}

\subsection{Постановка задачи}

По столбцу \(X_4\) проверяется гипотеза согласия с экспоненциальным законом распределения.

\subsubsection{Гипотезы}
\begin{itemize}
    \item \(H_0: X_4 \sim \text{Exp}(\lambda = 0.106)\) — выборка взята из экспоненциального распределения с параметром \(\lambda = 0.106\) (параметр задан в карточке варианта).
    \item \(H_1:\) распределение \(X_4\) не является экспоненциальным с \(\lambda = 0.106\).
\end{itemize}

\subsubsection{Уровень значимости}
\(\alpha = 0.05\).

\subsection{Выбор критерия}

\textbf{Выбранный критерий:} критерий согласия Пирсона (\(\chi^2\)) (приложение A.8), реализованный функцией \texttt{scipy.stats.chisquare}.

\subsection{Параметр распределения}

Параметр \(\lambda = 0.106\) \textbf{задан} в карточке варианта и не оценивается по выборке.

\subsection{Разбиение на интервалы}

Число интервалов выбрано по правилу Стерджеса:
\[
k = 1 + \lfloor \log_2 n \rfloor = 1 + \lfloor \log_2 106 \rfloor = 7.
\]

Используем равновероятностное разбиение: границы интервалов выбраны так, чтобы при справедливости \(H_0\) вероятность попадания в каждый интервал была одинакова (\(p_k = 1/7\)). Границы интервалов определяются как квантили экспоненциального распределения с \(\lambda = 0.106\):
\[
a_k = F^{-1}_{\text{Exp}}\left(\frac{k}{7}\right) = -\frac{1}{0.106} \ln\left(1 - \frac{k}{7}\right), \quad k = 0, 1, \dots, 7.
\]

\subsection{Таблица частот}

\begin{table}[H]
\centering
\caption{Наблюдаемые и ожидаемые частоты для \(X_4\)}
\label{tab:chi2}
\begin{tabular}{ccc}
\toprule
Интервал & \(n_k\) (набл.) & \(np_k\) (ожид.) \\
\midrule
\([0.00, 1.45)\) & 21 & 15.14 \\
\([1.45, 3.17)\) & 17 & 15.14 \\
\([3.17, 5.28)\) & 18 & 15.14 \\
\([5.28, 7.99)\) & 11 & 15.14 \\
\([7.99, 11.82)\) & 20 & 15.14 \\
\([11.82, 18.36)\) & 13 & 15.14 \\
\([18.36, +\infty)\) & 6 & 15.14 \\
\midrule
Итого & 106 & 106.00 \\
\bottomrule
\end{tabular}
\end{table}

Все ожидаемые частоты \(np_k = 15.14 \geq 5\), поэтому объединение интервалов не требуется.

\subsection{Вычисление статистики критерия}

Согласно приложению A.8, статистика критерия Пирсона имеет вид:
\[
\chi^2 = \sum_{k=1}^{m} \frac{(n_k - np_k)^2}{np_k} \approx 11.547.
\]

Число степеней свободы (приложение A.8):
\[
\nu = m - 1 = 7 - 1 = 6,
\]
где \(m = 7\) — число интервалов. Поскольку параметр \(\lambda = 0.106\) задан заранее, коррекция на число оценённых параметров не требуется.

\subsection{Принятие решения}

Согласно приложению A.3, используем p-value:
\[
p\text{-value} = P(\chi^2_6 \geq 11.547) \approx 0.073.
\]

\[
p\text{-value} = 0.073 > \alpha = 0.05.
\]

Следовательно, нет оснований отвергнуть \(H_0\).

\subsection{Вывод}

\textbf{Нет оснований отвергнуть нулевую гипотезу} \(H_0\) на уровне значимости \(\alpha = 0.05\). Данные \(X_4\) согласуются с экспоненциальным распределением \(\text{Exp}(0.106)\). Следует отметить, что p-value близко к пороговому значению (\(p = 0.073\)); наибольший вклад в статистику вносит последний интервал \([18.36, +\infty)\), в котором наблюдается лишь 6 значений при ожидаемых \(15.14\).

\subsection{Ошибка первого рода}

В данной задаче ошибка первого рода (приложение A.2) означает: \textbf{отвергнуть гипотезу о согласии \(X_4\) с экспоненциальным распределением \(\text{Exp}(0.106)\), когда на самом деле данные подчиняются этому закону}. Вероятность такой ошибки \(\alpha = 0.05\).

\section{Итоговый вывод}

В рамках расчётно-графической работы №2 по проверке статистических гипотез (вариант A-5, \(n = 106\), \(\alpha = 0.05\)) было выполнено четыре задания:

\begin{enumerate}
    \item \textbf{Параметрическая проверка гипотезы о равенстве математических ожиданий \(X_1\) и \(X_2\):}
    использован двухвыборочный t-критерий Стьюдента (приложение A.4, функция \texttt{scipy.stats.ttest\_ind}).
    Гипотеза \(H_0: \mathbb{E}X_1 = \mathbb{E}X_2\) \textbf{не отвергнута} (\(h = -1.544\), \(p = 0.124\)).
    Статистически значимых различий между средними не обнаружено.

    \item \textbf{Проверка гипотезы о математическом ожидании \(X_3\):}
    использован одновыборочный t-критерий Стьюдента (приложение A.5, функция \texttt{scipy.stats.ttest\_1samp}).
    Гипотеза \(H_0: \mu = 75.24\) \textbf{не отвергнута} (\(h = -0.502\), \(p = 0.616\)).
    Данные хорошо согласуются с предположением \(\mu = 75.24\).

    \item \textbf{Непараметрический критерий для \(X_1\) и \(X_2\):}
    использован U-критерий Манна–Уитни (приложение A.7, функция \texttt{scipy.stats.mannwhitneyu}).
    Гипотеза \(H_0\) об одинаковости распределений \textbf{не отвергнута} (\(p = 0.155\)).
    Выводы параметрического и непараметрического подходов совпали, что свидетельствует об устойчивости результата.

    \item \textbf{Критерий согласия для \(X_4\):}
    использован критерий \(\chi^2\) Пирсона (приложение A.8, функция \texttt{scipy.stats.chisquare}).
    Проверялась гипотеза \(H_0: X_4 \sim \text{Exp}(\lambda = 0.106)\), параметр задан.
    Гипотеза \textbf{не отвергнута} (\(\chi^2 = 11.547\), \(\nu = 6\), \(p = 0.073\)).
    Экспоненциальная модель с заданным параметром не противоречит данным, хотя p-value близко к границе принятия решения.

\end{enumerate}

\textbf{Общий итог:} во всех четырёх проверках нулевые гипотезы не были отвергнуты на уровне значимости \(\alpha = 0.05\). Это означает, что:
\begin{itemize}
    \item Различия между \(X_1\) и \(X_2\) не являются статистически значимыми — их можно рассматривать как выборки из распределений с одинаковыми средними. Оба критерия (параметрический и непараметрический) дали согласованные результаты.
    \item Данные \(X_3\) не противоречат предположению о среднем значении \(\mu = 75.24\).
    \item Экспоненциальное распределение с параметром \(\lambda = 0.106\) (заданным в варианте) в целом согласуется с \(X_4\), однако значение p-value (\(0.073\)) относительно близко к порогу — при более строгом уровне значимости (\(\alpha = 0.10\)) гипотеза была бы отвергнута.
\end{itemize}

Все критерии применены в соответствии с приложениями A.2–A.8: выполнена проверка предположений, вычислены наблюдаемые значения статистик и p-value, сформулированы содержательные выводы с интерпретацией возможной ошибки первого рода.
